%%%%%%%%%%%%%%%%%%%%%%%%%%%%%%%%%%%%%%%%%%%%%%%%%%%%%%%%%%%%%%%%%%%%%%%%%%%%%%%%
%%  content/01-introduction.tex — فصل ۱: مقدمه و بیان مسئله
%%
%%  هر \section یک اسلاید جداکننده می‌سازد (با sectionpage=false خاموش می‌شود)
%%  و نام همان فصل در نوار بالای عنوان اسلایدها تکرار می‌شود.
%%%%%%%%%%%%%%%%%%%%%%%%%%%%%%%%%%%%%%%%%%%%%%%%%%%%%%%%%%%%%%%%%%%%%%%%%%%%%%%%
\section{مقدمه و بیان مسئله}

%%%%%%%%%%%%%%%%%%%%%%%%%%%%%%%%%%%%%%%%%%%%%%%%%%%%%%%%%%%%%%%%%%%%%%%%%%%%%%%%
\begin{frame}{صورت مسئله}{از فهرست پیوندها تا پاسخ}
  \begin{itemize}
    \item کاربر پرسشی به \term{زبان طبیعی} می‌پرسد، اما موتور جست‌وجو فهرستی از
          پیوندها بازمی‌گرداند؛ یافتن پاسخ در میان آن‌ها بر دوش خود کاربر است.
    \item مدل‌های زبانی بزرگ پاسخ روان می‌دهند، ولی دانش‌شان به \hl{زمان آموزش}
          محدود است و در نتیجه از رخدادها و مستندات تازه بی‌خبرند.
    \item همین محدودیت به \term{توهم‌زایی} می‌انجامد: پاسخی که درست به نظر می‌رسد
          اما پشتوانهٔ قابل بررسی ندارد.
  \end{itemize}

  \begin{PTcallout}
    پرسش پژوهش: چگونه می‌توان \hl{روانی} پاسخ مدل زبانی را با \hl{تازگی} و
    \hl{قابلیت استناد} وب یکی کرد؟
  \end{PTcallout}
\end{frame}

%%%%%%%%%%%%%%%%%%%%%%%%%%%%%%%%%%%%%%%%%%%%%%%%%%%%%%%%%%%%%%%%%%%%%%%%%%%%%%%%
\begin{frame}{کاستی سامانه‌های موجود}
  \begin{itemize}
    \item \term{موتورهای کلیدواژه‌ای}: بازیابی سریع، اما ناتوان از درک منظور
          پرسش‌های چندبخشی.
    \item \term{دستیارهای گفت‌وگویی بستهٔ وب}: پاسخ یکدست، بدون ارجاع بررسی‌پذیر
          به سند اصلی.
    \item \term{سامانه‌های تجاری پرسش و پاسخ}: جعبهٔ سیاه‌اند؛ نه سنجش‌پذیرند و نه
          برای زبان و دامنهٔ خاص قابل تنظیم.
  \end{itemize}

  \begin{PTcons}
    \begin{itemize}
      \item نبود سنجهٔ روشن برای اینکه هر جزء پاسخ از کدام سند آمده است.
      \item هزینهٔ بالای فراخوانی مدل، وقتی همهٔ اسناد بازیابی‌شده به آن سپرده شوند.
    \end{itemize}
  \end{PTcons}
\end{frame}

%%%%%%%%%%%%%%%%%%%%%%%%%%%%%%%%%%%%%%%%%%%%%%%%%%%%%%%%%%%%%%%%%%%%%%%%%%%%%%%%
\begin{frame}{هدف و مشارکت‌های این پژوهش}
  \begin{enumerate}
    \item طراحی سامانه‌ای مبتنی بر وب که پرسش زبان طبیعی را می‌گیرد و
          \hl{پاسخ مستند} با ارجاع به منبع بازمی‌گرداند.
    \item زنجیرهٔ \term{بازیابی، بازرتبه‌بندی و تولید} با اجزای جایگزین‌پذیر، تا
          هر مرحله جداگانه سنجیده شود.
    \item \term{بازرتبه‌بندی دومرحله‌ای} برای کوتاه کردن بافتار ورودی مدل و کاهش
          هزینهٔ هر پرسش.
    \item ارزیابی مقایسه‌ای روی مجموعه‌ای از \hl{۴۰ پرسش} در برابر یک سامانهٔ
          تجاری مرجع.
  \end{enumerate}

  \begin{PTresult}
    دقت پاسخ‌دهی سامانه \hl{۹۵٪} در برابر \hl{۷۷٫۵٪} سامانهٔ مرجع، با بافتاری
    کوتاه‌تر و هزینهٔ کمتر.
  \end{PTresult}
\end{frame}
